% Direct-authority diplomatic French for printed p.182.
\SourcePage{182}

\LayerTarget{tohoku:proof:proposition-3.10.1:p2}
$H^p(\mathop{\underset{\to}{\lim}} F_iC(A))$, ils sont donc isomorphes puisque le foncteur
$\mathop{\underset{\to}{\lim}}$ sur la catégorie des systèmes inductifs sur $I$ à valeurs dans
$\mathbf C'$ est exact (prop. 1.8) et en particulier permute à la formation
de homologies de complexes.

\LayerTarget{tohoku:corollary:3.10.1:1}
\textsc{Corollaire 1.} \textit{Soit $X$ un espace muni d'une famille
paracompactifiante $\Phi$. On a alors, pour tout faisceau abélien $F$ sur
$X$ :}
\LayerTarget{tohoku:eq:corollary-3.10.1-1:support-limit}
\[
H^p_{\Phi}(X,F)=\varinjlim\limits_U H^p(X,F_U)
\]
\textit{la limite inductive étant prise suivant l'ordonné filtrant des ouverts
$U$ de $X$ dont l'adhérence est dans $\Phi$. \textup{(}$F_U$ désigne le faisceau sur
$X$ dont la restriction à $U$ est $F|U$, et la restriction à $\mathrm C U$ est \textup{0)}.}

\LayerTarget{tohoku:proof:corollary-3.10.1-1}
En vertu du théorème 3.5.1. on a
$H^p(X,F_U)=H^p_{\Phi_U}(U,F)$, où $\Phi_U$ est l'ensemble des parties de $U$
fermées dans $X$. Posant $\Gamma_{\Phi_U}(F)=H^0(X,F_U)$, on peut donc aussi
écrire $H^p(X,F_U)=R^p\Gamma_{\Phi_U}(F)$ (compte tenu de proposition
3.1.3), d'oû en vertu de prop. 3.10.1 :
$\mathop{\underset{\to}{\lim}} H^p(X,F_U)=R^p\Gamma_{\Phi}(F)$ puisque
$\mathop{\underset{\to}{\lim}}\Gamma_{\Phi_U}(F)=\Gamma_{\Phi}(F)$, cqfd.
\LayerTarget{tohoku:remark:corollary-3.10.1-1:cartan}
--- Le corollaire précédent est parfois utile pour ramener la cohomologie
``à supports dans $\Phi$'' à la cohomologie à supports quelconques, et m'avait
été signalé par M. Cartan.

\LayerTarget{tohoku:corollary:3.10.1:2}
\textsc{Corollaire 2.} \textit{Soient $X$ un espace topologique, $Y$
une partie de $X$ admettant un système fondamental de voisinages paracompacts
\textup{(}il suffit par exemple que $X$ soit métrisable, ou localement compact et
paracompact\textup{)}. Alors pour tout faisceau abélien $F$ sur $X$, on a}
\LayerTarget{tohoku:eq:corollary-3.10.1-2:neighborhood-limit}
\[
H^p(Y,F)=\mathop{\underset{\to}{\lim}} H^p(U,F)
\]
\textit{la limite étant prise suivant l'ensemble ordonné filtrant décroissant
des voisinages ouverts de $U$ dans $X$.}

\LayerTarget{tohoku:proof:corollary-3.10.1-2}
En effet, il résulte de l'hypothèse que
$H^0(Y,F)=\mathop{\underset{\to}{\lim}} H^0(U,F)$ ([11], prop. 2.2.1), or les foncteurs dérivés
de $F\to H^0(U,F)$ sont les $H^p(U,F)$, de sorte que le corollaire 2
est un cas particulier de la proposition. On notera que l'on a aussi
$H^0(Y,F)=\mathop{\underset{\to}{\lim}} H^0(U,F)$, et par suite la validité du corollaire 2, si
$Y$ est \textit{fermé} et admet \textit{un} voisinage paracompact
(démonstration analogue à celle de [11]) ; dans loc. cité, on trouvera aussi
un contre-exemple simple (avec $p=0$) relatif au cas où on ne fait pas
d'hypothèse de paracompacité.

\LayerTarget{tohoku:result:3.10.2:preliminary}
A titre de complément, indiquons encore sans démonstration le résultat
suivant, cas particulier de résultats généraux sur les systèmes projectifs.
Soit $X$ un espace localement compact, considérons l'ensemble filtrant
croissant des parties ouvertes relativement compactes $U$ de $X$, alors pour
tout faisceau abélien $F$ sur $X$, les homomorphismes de restriction
$H^p(X,F)\to H^p(U,F)$ définissent des homomorphismes canoniques
(qui sont évidemment des homomorphismes de $\partial$-foncteurs) :
\LayerTarget{tohoku:eq:3.10.2}
\[
\tag{3.10.2}
H^p(X,F)\to\mathop{\underset{\leftarrow}{\lim}} H^p(U,F)
\]
qui sont évide nment bijectifs pour $p=0$.

\TohokuContinuation{183}
